\section{Executor independence and the justification boundary}\label{sec:executor-independence}

The mechanics in this paper are not defined by the use of a large language model. They are conditions on scientific state transitions. We therefore distinguish the \emph{mechanic} from the \emph{executor} that proposes or computes a transition. An executor may be a human researcher, a deterministic program, a theorem prover, a learned model, a large language model (LLM), or a hybrid workflow. This distinction matters because the same epistemic error can arise under any executor: provenance can be lost, incompatible contexts can be merged, an obstruction can be discarded, or an unsupported proposal can be mistaken for an authority-bearing update.

Let a mechanic be a typed partial transition
\[
M:(S,I)\rightharpoonup(S',O,R),
\]
where $S$ is research state, $I$ is an input object, $O$ is an output, and $R$ is a receipt containing the declared evidence and transition provenance. Different implementations may realize the same abstract mechanic,
\[
M_{\rm human},\qquad M_{\rm symbolic},\qquad M_{\rm prover},\qquad M_{\rm LLM},\qquad M_{\rm hybrid}.
\]
Implementation identity belongs in provenance and cost accounting, but it does not by itself change the scientific-authority coordinates. The non-escalation results in this paper therefore apply to classes of transitions rather than to a model family. If an operator has no certified authority effect, finite compositions of that operator class leave the canonical scientific projection unchanged.

This framing has two practical consequences. First, the framework admits fully non-neural executions. The current reference system already contains LLM-free paths for finite-basis saturation, support-hypergraph traversal, obstruction checking, certificate validation and theorem-kernel discharge. Classical knowledge-representation systems offer further algorithmic parents: assumption-based truth maintenance maintains multiple assumption environments rather than collapsing inconsistent hypotheses into one world \cite{dekleer1986}; abstract argumentation separates attacks, defence and acceptability \cite{dung1995}; abstract interpretation reasons over sound approximations and fixpoints \cite{cousot1977}; and proof assistants such as Lean separate extensible automation from a small proof-checking kernel \cite{demoura2021lean}. These examples motivate a useful engineering rule for Orion: use an LLM only where a learned generative adapter is actually needed, and prefer a deterministic checker when the relevant relation is already formal.

Second, LLM output is treated as a proposal surface, not a privileged epistemic channel. This is compatible with recent scientific-agent systems in which language models coordinate search, tool use or experimental planning \cite{boiko2023,bran2024,gottweis2026coscientist}. Such systems demonstrate that LLMs can be productive research executors, but they do not eliminate the distinction between generation and justification. In Orion, an LLM may propose a claim, a representation, an analogy, a search action or a mechanic revision. The proposal is then subject to the same provenance, compatibility, evidence and authority rules that would apply if the proposal came from a human or a symbolic program.

\subsection{Source, familiarity and repeated exposure}

Executor independence also clarifies a failure mode that is easy to misread as specifically neural. Repetition can make a statement familiar without adding an independent evidence root. Source-monitoring research shows that judgments about the origin of remembered information are reconstructive and can misattribute familiarity \cite{johnson1993source}. Work on learning from testimony likewise shows that source competence and reliability are task- and context-dependent rather than interchangeable with the content of a claim \cite{harris2018testimony}. We therefore treat retrieval count, repetition and source familiarity as authority-inert metadata. A source may be a useful routing prior, but repeated exposure is not new evidence.

The first Paper-I repetition-attack experiment tests this separation rather than assuming it. Under its frozen attack model, ten submitted identifiers collapsed to only eight normalized identities, giving a repetition ratio of $0.2$ against a registered $0.5$ gate. The terminal is therefore negative. That result is retained as a negative epoch and motivates a materially different identity-and-lineage mechanism; it is not evidence that the source-monitoring firewall is already solved.

\subsection{Decision-relative views and bounded stopping}

The executor-independent view also narrows two other claims. First, a task-conditioned projection is justified only relative to a declared decision or query class. Blackwell's comparison of experiments provides the classical decision-theoretic precedent for comparing information structures by the decisions they support rather than by raw descriptive richness \cite{blackwell1953}. Paper~I's projection-collision theorem is the local obstruction: if two worlds are identified by a projection but require different registered actions, that projection is insufficient for the task.

Second, saturation is necessarily scoped. A finite transcript cannot establish that no undiscovered mechanism exists in an unrestricted open world. The strongest legitimate stopping claim is therefore relative to a declared search/action basis, budget and expiry condition. When statistical successors are launched adaptively, programme-level evidence also needs to account for optional continuation; anytime-valid inference provides one principled family of tools for doing so \cite{ramdas2023}. These qualifications do not weaken the framework. They define the exact object that can be checked by humans, algorithms or learned executors alike.
